\section{Verifier Interface and Dispatch}
\label{sec:verifier-interface}

\subsection{The \texttt{P3QVerifier} Go interface}

The unification of the three kinds onto one slot is mediated by
a Go-side interface at
\texttt{lux/standard/precompiles/p3q\_family/}. The interface is
the single audit boundary for adding a kind: implement the four
methods, register at \texttt{init()}, done.

\begin{lstlisting}[language={}]
// P3QVerifier is the per-kind verifier abstraction.
//
// Each P3QVerifier implementation is registered at process boot
// into a kind-byte keyed dispatch table loaded by the single
// P3Q precompile at slot 0x012205.
type P3QVerifier interface {
    // Verify returns true iff sig is a valid threshold signature
    // from groupPubKey on messageHash. The implementation must
    // be constant-time over secret inputs (there are none; both
    // sig and pk are PUBLIC calldata) and must not panic on
    // malformed input -- malformed inputs return false.
    Verify(sig []byte, messageHash [32]byte, groupPubKey []byte) bool

    // GasCost returns the per-call precompile gas charge for
    // this kind. The P3Q dispatcher charges this amount BEFORE
    // calling Verify, so an out-of-gas caller never observes the
    // verifier's behavior on its specific input.
    GasCost() uint64

    // Kind returns the calldata discriminator byte:
    //   0x01 = Pulsar     (Module-LWE)
    //   0x02 = Corona     (Module-LWE)
    //   0x03 = Magnetar   (hash-based)
    Kind() byte

    // PrimitiveName returns "pulsar", "corona", or "magnetar".
    // Used for metric labels, log lines, and audit trails.
    PrimitiveName() string
}
\end{lstlisting}

The interface is minimal by design: four methods, three of which
are constant-returning getters. The verifier core lives in the
\texttt{Verify} method; the gas calibration lives in
\texttt{GasCost}; the dispatch discriminator lives in
\texttt{Kind}. There is no shared base class, no inheritance, no
mutable verifier state -- each verifier is value-typed and
independently complete (composition over inheritance per
LP-200).

\subsection{Dispatch flow}

When the \texttt{0x012205} precompile is called:

\begin{enumerate}[leftmargin=2em,nosep]
\item Read the leading \texttt{kind} byte from calldata
  (offset~0).
\item Look up the registered \texttt{P3QVerifier} for that kind
  in the kind-keyed dispatch table.
\item Charge \texttt{Verifier.GasCost()} (per-kind, returned by
  the impl). The uniform per-byte adder is added separately by
  the precompile's \texttt{RequiredGas}.
\item Parse the wire format down to
  \texttt{(sig, messageHash, groupPubKey)} using the layout
  appropriate to the kind (the per-kind size constants are
  baked into the implementation; the framing -- BE32 length
  fields, trailing 32-byte messageHash -- is uniform).
\item Call \texttt{Verifier.Verify(sig, messageHash, pk)}.
\item Return the boolean result encoded as a 32-byte EVM-ABI
  \texttt{bool} (\texttt{abiTrue} or \texttt{abiFalse}).
\end{enumerate}

This is the single point of truth for the family. Three verifier
implementations, one slot, one dispatch path, one ABI shape, one
audit surface. Adding a fourth PQ-family kind (a hypothetical
code-based threshold via HQC-Sig variants, for example) is one
new \texttt{P3QVerifier} impl and one new \texttt{kind} byte
registered in the same dispatch table; no new precompile slot,
no new Solidity function shape, no new audit boundary at the
EVM-side.

\subsection{Per-kind gas table}

\begin{table}[h]
\centering
\small
\begin{tabular}{@{}lrrr@{}}
\toprule
\textbf{kind} & \textbf{Name} & \textbf{Gas charge} & \textbf{Verify CPU} \\
\midrule
$0x01$ & Pulsar   & \SI{50000}{}    & $\sim$\SI{80}{\micro\second} \\
$0x02$ & Corona   & \SI{5000000}{}  & $\sim$\SI{1.6}{\milli\second} \\
$0x03$ & Magnetar & \SI{30000}{}    & $\sim$\SI{1.92}{\milli\second} \\
\bottomrule
\end{tabular}
\caption{Per-kind gas table. Corona's \SI{5}{M} charge is the
$\sim$$100\times$ multiplier over Pulsar reflecting the CPU
verify-cost ratio; Magnetar's \SI{30000}{} reflects the fact
that the SLH-DSA hash-tree walk is memory-bandwidth-bound and
parallelizes harder on modern CPUs than the lattice NTT
inverse.}
\label{tab:gas-per-kind}
\end{table}

The asymmetry between Corona's \SI{5}{M} and Magnetar's
\SI{30000}{} -- despite comparable CPU verify times -- reflects
two distinct calibrations:

\begin{itemize}[leftmargin=2em,nosep]
\item Corona's verify dominates a single core because the NTT
  inverse + polynomial multiplication mod $q$ over the larger
  Ring-LWE coefficients is a tight cache-resident loop with
  limited ILP. \SI{1.6}{\milli\second} on a single core is
  $\sim$$141\times$ the Pulsar verify; the gas charge rounds
  to $\sim$$100\times$.
\item Magnetar's verify is a hash-tree walk: hash function
  evaluation is highly parallelizable, well-vectorized by
  modern CPUs, and memory-bandwidth-bound. \SI{1.92}{\milli\second}
  wall-clock is dominated by sequential dependencies in the
  WOTS+ structure but the per-operation cost is lower than the
  lattice multiplication. The gas charge reflects the modest
  per-call CPU cost.
\end{itemize}

This is the same calibration discipline LP-218 uses for Pulsar:
gas reflects CPU-class verifier cost, not GPU-class. Operators
with all-Blackwell validator sets may petition for discounted
gas in a future amendment; LP-220 does not codify it.

\subsection{Validator registration}

Each P3Q-conformant validator MUST register a
\texttt{P3QVerifier} for every kind it will accept on-chain.
Misconfiguration -- a validator missing the Corona verifier
when Z-Chain receives a rollup-batch tx that calls
\texttt{P3Q.verify(0x02, ...)} -- is a configuration error, not
a soundness error: the validator's EVM cannot validate the
block, and consensus excludes that validator's vote on the
block. LP-202 tier-degradation handles the recovery (the
validator votes \texttt{nil} on the block; consensus continues
with the remaining validators).

The reference \texttt{p3q.Run} ships KindPulsar wired and
KindCorona / KindMagnetar reserved-but-not-wired; this is the
genesis state. As the Corona and Magnetar verifier libraries
land in subsequent versions of \texttt{luxfi/corona} and
\texttt{luxfi/magnetar}, their P3Q dispatch entries are wired
in, and the wired-state advances. The wire format and the
Solidity ABI are stable across the wired/not-wired transition;
only the verify behavior changes (from \texttt{false} to the
actual cryptographic outcome).

\subsection{Failure modes}

\begin{table}[h]
\centering
\small
\begin{tabular}{@{}p{4.8cm}p{4cm}p{4.5cm}@{}}
\toprule
\textbf{Failure} & \textbf{Trigger} & \textbf{Effect} \\
\midrule
Unknown \texttt{kind} byte & caller passes
  \texttt{kind} not in $\{0x01, 0x02, 0x03\}$ &
  precompile returns \texttt{abiFalse}; caller observes \texttt{false} \\
Reserved-but-not-wired kind & caller passes \texttt
  or \texttt before verifier lands &
  precompile returns \texttt{abiFalse}; caller observes \texttt{false} \\
Malformed wire layout & caller passes truncated/oversized
  fields & precompile returns \texttt{abiFalse} \\
Sig-length mismatch (Pulsar) & \texttt{sigLen} $\neq$ FIPS-204
  expected size for the mode byte & precompile returns
  \texttt{abiFalse} \\
FIPS verifier reject & sig genuinely does not verify under the
  group pubkey & precompile returns \texttt{abiFalse} \\
Out-of-gas & caller supplies
  \texttt{gas} $<$ \texttt{RequiredGas(input)} & precompile
  returns \texttt{contract.ErrOutOfGas}; caller transaction reverts \\
Group pubkey rotation skew & sequencer rotated keys without
  updating the on-chain group pubkey & precompile returns
  \texttt{abiFalse} until rotation tx lands \\
\bottomrule
\end{tabular}
\caption{Failure modes. Every cryptographic failure is
\texttt{abiFalse}; only out-of-gas reverts. This is the LP-218
``no revert on cryptographic failure'' contract -- the
Solidity caller decides whether \texttt{false} is fatal to its
batch.}
\label{tab:failures}
\end{table}
